\section{Introduction}
\label{sec:introduction}

Let $m$ be normalized Lebesgue area on $\T^2=\R^2/\Z^2$.
The standard sine map is the area-preserving diffeomorphism
\begin{equation}\label{main:map}
 f_k(x,y)=\bigl(x+y+k\sin(2\pi x),\;y+k\sin(2\pi x)\bigr)
 \pmod{\Z^2},\qquad k\in\R.
\end{equation}
We consider its Kolmogorov--Sinai entropy $h_m(f_k)$ for this fixed
invariant measure. The distinction between topological complexity and
chaos visible to area is essential: a hyperbolic set may have zero area,
and even full Hausdorff dimension does not imply positive area.

\begin{theorem}\label{main:entropy}
There exists $k_0>0$ such that
\[
 h_m(f_k)>0\qquad\text{for every }k\ge k_0.
\]
Equivalently, for every such $k$ the largest Lyapunov exponent of $f_k$
is positive on a set of positive $m$-measure.
\end{theorem}

In particular the parameter set with positive metric entropy contains
an interval of positive length. Thus \cref{main:entropy} resolves
affirmatively Sinai's positive metric entropy conjecture in the
formulation of Berger and Turaev
\cite[Conjecture~0.4]{BergerTuraev2019}, and gives the stronger
all-sufficiently-large-parameter conclusion. This large-parameter
formulation is also recorded by Obata \cite{Obata2021}.
Berger and Turaev's recurrence
$(x,y)\mapsto(2x-y+a\sin(2\pi x),x)$ is conjugate to
\cref{main:map} by $(x,y)\mapsto(x,x-y)$, with $a=k$.
The assertion concerns both parameter measure and, at each parameter,
phase-space area. It makes no claim of ergodicity, almost-everywhere
positive exponents, or a uniform positive lower bound on entropy.

\subsection*{Context}

The standard map arose as a model of conservative instability
\cite{Chirikov1979}. Its large-parameter regime combines strong local
expansion with recurrent losses of derivative growth. Duarte's
large-parameter analysis constructs hyperbolic basic sets and exhibits
elliptic-island and tangency phenomena in the standard family
\cite{Duarte1994}. Gorodetski \cite{Gorodetski2012} obtained transitive
invariant sets of full Hausdorff dimension for a residual set of large
parameters. These results reveal extensive deterministic structure,
but neither hyperbolic sets nor full dimension imply positive area;
residual parameter sets also need not have positive Lebesgue measure.

Other entropy results use different measures or notions. Obata
\cite{Obata2021} established uniqueness of the measure of maximal
entropy for sufficiently large standard-map parameters, without
identifying that measure with area. Oliveira
\cite[Theorem~5.1]{Oliveira2024} proved positive topological entropy
for every $k>0$ with $k\ne2/\pi$ in the normalization of
\cref{main:map}. Positive topological entropy does not by itself give positive
entropy for the specific invariant measure $m$. Berger and Turaev
\cite{BergerTuraev2019} proved Herman's positive-entropy conjecture by
conservative perturbations, including perturbations near standard
maps; these are not restricted to amplitude changes in \cref{main:map}.
Related hyperbolicity is available for Anosov-driven skew products
with standard-map fibers
\cite{BergerCarrasco2014} and randomly perturbed standard maps
\cite{BlumenthalXueYoung2017}. Coupling, noise, and smooth
perturbation change the dynamical system or its probability law. The
question here concerns Lebesgue area for the exact, fixed sine map at
each sufficiently large amplitude.

We use the established entropy formula, not a new entropy theory.
For a smooth area-preserving diffeomorphism of the compact torus,
the multiplicative ergodic theorem \cite{Oseledets1968} gives
almost-everywhere Lyapunov exponents. The determinant-one condition
makes them $\lambda_+(z)$ and $-\lambda_+(z)$, with
$\lambda_+(z)\ge0$. The entropy formula of Pesin, in the absolutely
continuous invariant-measure form
\cite{Pesin1977,Mane1981,Mane1983Errata}, gives
\begin{equation}\label{main:pesin}
 h_m(f_k)=\int_{\T^2}\lambda_+(z)\dd m(z).
\end{equation}
The map is $C^\infty$, the manifold is compact without boundary,
and $m$ is an absolutely continuous invariant probability. The
derivative and inverse-derivative logarithms are bounded at each fixed
parameter, so all integrability hypotheses hold. Neither ergodicity
nor a prior assumption of nonzero exponents is required. Consequently
the task is to prove positivity of $\lambda_+$ on positive area.

\begin{corollary}[Nonuniformly hyperbolic ergodic component]
\label{main:components}
For each $k\ge k_0$ there is a measurable $f_k$-invariant set $E=E(k)$
with $m(E)>0$ such that $(E,m(E)^{-1}m|_E,f_k)$ is ergodic and $f_k$ has
one positive and one negative Lyapunov exponent almost everywhere on $E$.
For some integer $N=N(k,E)\ge1$, there is a measurable partition
$E=E_0\sqcup\cdots\sqcup E_{N-1}$ such that
$f_k(E_j)=E_{j+1\bmod N}$ and
$(E_j,m(E_j)^{-1}m|_{E_j},f_k^N)$ is Bernoulli for every $j$.
All invariance and partition statements are modulo $m$-null sets;
each $E_j$ has area $m(E)/N>0$.
\end{corollary}

\begin{proof}
Fix $k\ge k_0$ and let $H_k$ be the invariant Oseledets set where
$\lambda_+>0$. It has positive area by \cref{main:entropy}, while
$\det Df_k=1$ makes the other exponent $-\lambda_+<0$ there.
The normalized restriction $\nu_k=m(H_k)^{-1}m|_{H_k}$ is an absolutely
continuous invariant probability on the compact torus, with no zero
Lyapunov exponents. Pesin's ergodic and Bernoulli decomposition theorems,
as stated in \cite[Theorems~3.3--3.4]{KatokBurns1994}, apply to $\nu_k$
and give a positive-$\nu_k$ ergodic component $E$ and its finite cyclic
partition. Since $\nu_k(E)=m(E)/m(H_k)>0$, the component has positive
area; normalizing $\nu_k$ on $E$ or on a partition piece agrees with
normalizing $m$. Area preservation gives $m(E_j)=m(E)/N$.
\end{proof}

\subsection*{Proof mechanism}

The proof has two quantitative inputs and a multiscale conclusion.
In transfer coordinates the derivative is a product of determinant-one
matrices $T_{i,j}$. With the one-step bound $M=2\pi k+4$, put
$g(i,j)=\log_M\norm{T_{i,j}}$. The bound $0\le g(i,j)\le|i-j|$
and the triangle inequality make $g$ a stationary semimetric on orbit
time. We track the expected shortfall $1-\E g(0,n)/n$ from maximal
growth at each length $n$.

First, \cref{can:theorem} bounds the probability that two long segments
grow nearly at the maximal rate but their concatenation almost
completely cancels. A cancellation supplies a linearized solution that
is small at both endpoints. Matched expansion scales on its two sides
give many local tests. A uniform analytic scarcity estimate, including
the degenerate cases of nearly cancelling reciprocal cosines, converts
those tests into exponential probability decay. Finite-time direction
and tangency geometry in the standard family has precedents in
\cite{BloorLuzzatto2009}; the matched-scale scarcity estimate used
here is proved below, rather than imported from that work.

Second, \cref{bridge:theorem} gives a weighted product estimate for
observations near opposite ends of a rapidly growing segment. Finite
Dirichlet problems construct two transverse contracting graph families.
Their uniformly controlled slice Jacobians yield the estimate by
changes of variables. This step does not presume mixing or independence
of orbit segments. For comparison, Blumenthal
\cite{Blumenthal2021} obtains fixed-parameter correlation
bounds that are useful over a finite parameter-dependent time range.
Those bounds do not supply the weighted estimate across fast bridges in
\cref{bridge:product}, which applies at arbitrary sufficiently large
bridge lengths; its restriction to fast bridges is essential.

Suppose now that arbitrarily large parameters had zero exponent almost
everywhere. The shortfall at every fixed length tends to zero as the
parameter grows, whereas at each fixed zero-exponent parameter it tends
to one as the length grows. Cancellation scarcity bounds how much this
shortfall can jump when the length doubles. We can therefore select a
range of dyadic lengths over which it first increases and normalize by
the small terminal shortfall. Rescaling the distance arrays produces
measures that are locally finite away from the affine arrays
$d(s,t)=w|s-t|$, where every interval grows at one common linear rate.
The bridge estimate restricts which slow intervals can coexist with a
fast connecting segment. Together with cancellation scarcity and a
four-point inequality, it leaves three kinds of nonaffine limit arrays:
an entirely slow array, or one with one or two exterior unit-speed rays.

The final comparison uses a capped shortfall: it is linear with a small
slope near zero, and equals one on uniformly slow intervals. Sum its
whole-interval value minus the average of its two half-interval values
over the selected dyadic lengths. This telescopes to a positive
normalized terminal value, since the initial contribution vanishes.
The limiting sum decomposes into the backward dilates of the terminal
measure and a translation- and dilation-invariant remainder. The first
contribution is strictly less than that same terminal value: a nonzero
terminal measure retains a slow interval visible at arbitrarily fine
scales. The invariant remainder's contribution is nonpositive. For its two-ray shapes,
translation and dilation averaging compare fine-scale occupancy of the
slow interval with the smaller capped loss over long crossing intervals.
These two bounds contradict the positive telescoping balance.

Passing to the limit in that balance is itself a proof obligation.
Vague convergence away from affine arrays does not control weighted mass
approaching them. The cap is flat at slow affine speeds; the bridge
estimate removes intermediate affine speeds; and cancellation truncation
controls the remaining unit-speed array. Thus the same quantitative
inputs govern both the nonaffine shapes and the mass that might otherwise
escape them.

\Cref{sec:setup} fixes the transfer coordinates and growth conventions.
\Cref{can:section,loc:section} prove cancellation scarcity and its local
analytic estimate. \Cref{bridge:section} supplies the finite-graph
product estimate. \Cref{sec:scales} constructs the dyadic limit measures;
\cref{shape:section} classifies their support, and \cref{end:section}
completes the capped-deficit balance. These quantitative ingredients
are proved in the form needed here; no external standard-map entropy
assertion or independence hypothesis is imported.
